\section{讨论与结论}
\label{sec:conclusion}

LASA 的性能分布表明，片上物化重放有效支撑了高度异构的完整网络，但长规约层
仍决定主要执行时间。m13 与 m19 合计占 PL 时间 \LongLayerPLShare{}\%，未来
优化应优先减少这些层的反馈间隔和物化访问冲突。当前阵列原生覆盖标准
1$\times$1 和 3$\times$3 卷积；深度卷积、分组卷积和残差融合需要新的映射
方式。URAM 使用率达到 75\%，路由建立余量为 4 ps，因此进一步扩大缓存或阵列
需要同步考虑物理布局，而不是仅依赖算术资源余量。

从系统层面看，PL 已占 resident 时间的 96.2\%，A53 非卷积算子和后处理合计
只需 1.301 ms。这表明当前软硬件边界与算术密度相匹配：把不规则节点留在 A53
没有形成新的主瓶颈，继续优化应集中于长层卷积。当前量化参数与 FP32-416 的
精度差异描述模型状态，整数节点零失配则描述部署正确性；两种指标的分离使模型
参数能够迭代，同时保持相同硬件结构和评价方法。

本文围绕长 $K$ 卷积中“分段复用与逐轮片外搬运”的冲突提出 LASA。第一，
片上物化重放把 HWC tile 一次转换为阵列向量，使 DDR 输入流量从与
$P_lQ_l$ 相关降为与原始 tile 大小相关，并以层自适应分块覆盖完整13层。
第二，上下文调度通过权重提前装载、连续部分和反馈和直接切换缩短分段边界，
在三条正确性不变量下保持紧凑 HWC 输出顺序。

XCK26 上的 200 MHz 实现达到 \ResidentFPS{} FPS、\PLEffectiveGOPS{} GOPS 和
\ArrayUtilization{}\% 阵列利用率；\NodeRecords{} 条逐节点记录零失配，完整
5000 张原始检测头指标与主机参考一致。结果说明，一次物化、多轮片上重放和
连续上下文调度能够在固定规模阵列上支撑完整双尺度目标检测，并将理论数据复用
转化为稳定、可验证的端到端性能。
